% ---- score2 standard preamble (from: python3 score2_tools.py preamble) ----
\documentclass[10pt]{extarticle}
\usepackage[a4paper,margin=1.5cm]{geometry}
\usepackage[T1]{fontenc}
\usepackage{amsmath,amssymb}
\usepackage{graphicx}
\usepackage[export]{adjustbox}
\usepackage{multicol}
\usepackage{enumitem}
\usepackage{booktabs,array,tabularx}
\setlength{\columnsep}{0.7cm}
\setlength{\parindent}{0pt}
\setlength{\parskip}{3pt}
\pagestyle{empty}
\makeatletter
% figure registry (figdefs.tex is generated by the extract step)
\newcommand{\qfigdef}[4]{\@namedef{qf@file@#1}{#2}\@namedef{qf@w@#1}{#3}\@namedef{qf@wide@#1}{#4}}
\newcommand{\qf@check}[1]{\@ifundefined{qf@file@#1}{\PackageError{score2}{Unknown figure id #1}{}}{}}
\newcommand{\qf@img}[4][]{\includegraphics[width=#2\dimexpr\@nameuse{qf@w@#4}pt\relax,max width=#3,max height=0.45\textheight,#1]{\@nameuse{qf@file@#4}}}
% Full-width questions: finalize wraps every question that contains a wide figure in \qfullbegin...\qfullend,
% so the WHOLE question (stem, figures, options) is set across the page in reading order, then the two columns
% resume. The question's head - its number, stem and first figure (\qfullsplit marks the end of the head) -
% is kept on one page, so a number is never left at a page foot with its figure overleaf; the rest may break
% across pages normally, so no big gaps are left. Without a \qfullsplit, the whole question is kept together
% if it is under 45% of a page. Never write these yourself.
\newif\ifqf@full
\newif\ifqf@boxed
\newsavebox{\qf@box}
\newcommand{\qf@open}{\global\qf@fulltrue\par\global\qf@boxedtrue
  \setbox\qf@box\vbox\bgroup\hsize\textwidth\linewidth\textwidth}
\newcommand{\qf@emit}[1]{\par\egroup\global\qf@boxedfalse
  \ifdim\dimexpr\ht\qf@box+\dp\qf@box\relax<#1\textheight\noindent\usebox{\qf@box}\par\else\unvbox\qf@box\fi}
\newcommand{\qfullsplit}{\ifqf@boxed\qf@emit{0.85}\fi}
\newcommand{\qf@close}{\ifqf@boxed\qf@emit{0.45}\fi\global\qf@fullfalse}
\newcommand{\qfullbegin}{\end{multicols}\qf@open}
\newcommand{\qfullend}{\qf@close\begin{multicols}{2}}
\newcommand{\qfullstart}{\qf@open}
\newcommand{\qfulldone}{\qf@close}
% \qfig[scale]{id}: block figure, centred, directly where it is written (size set by the tool).
% Fallback if a wide figure is not inside a full-width question: break out of the columns in place - no floats.
\newcommand{\qfig}[2][1]{\qf@check{#2}%
  \ifqf@full\qf@blk{#1}{#2}\else
  \ifnum\@nameuse{qf@wide@#2}=1 \expandafter\@firstoftwo\else\expandafter\@secondoftwo\fi
  {\end{multicols}\begin{center}\qf@img{#1}{\textwidth}{#2}\end{center}\begin{multicols}{2}}%
  {\qf@blk{#1}{#2}}\fi}
\newcommand{\qf@blk}[2]{\par\smallskip{\centering\qf@img{#1}{\linewidth}{#2}\par}\smallskip}
% \qfigi[scale]{id}: inline figure (inside option lists / tables / minipages); never breaks columns.
\newcommand{\qfigi}[2][1]{\qf@check{#2}\qf@img[valign=c]{#1}{\linewidth}{#2}}
\makeatother
% structure macros - lint reads these to check completeness and figure placement
\newcommand{\papersec}[1]{\par\bigskip{\centering\Large\bfseries PAPER-#1\par}\medskip}
\newcommand{\partsec}[2]{\par\medskip\noindent{\large\bfseries PART-#1: #2}\par\smallskip}
\newcommand{\sect}[1]{\par\smallskip\noindent\textbf{#1}\par}
\newcommand{\qn}[1]{\par\smallskip\noindent\textbf{#1.}~}
\newcommand{\paraq}[2]{\par\smallskip\noindent\textbf{Paragraph for Questions #1 and #2}\par}
\newenvironment{qoptions}{\begin{itemize}[leftmargin=1.8em,itemsep=1pt,topsep=2pt,label={}]}{\end{itemize}}
\newcommand{\opt}[2]{\item[(#1)] #2}
% figure registry (generated by score2_tools.py finalize)
\qfigdef{p1_phy_s1_q02_1}{Score-II - Test 40 (16 Mar 2025)-img/p1_phy_s1_q02_1.png}{88.4}{0}
\qfigdef{p1_phy_s2_q05_1}{Score-II - Test 40 (16 Mar 2025)-img/p1_phy_s2_q05_1.png}{193.4}{0}
\qfigdef{p1_phy_s2_q06_2}{Score-II - Test 40 (16 Mar 2025)-img/p1_phy_s2_q06_2.png}{118.7}{0}
\qfigdef{p1_phy_s2_q06_1}{Score-II - Test 40 (16 Mar 2025)-img/p1_phy_s2_q06_1.png}{118.6}{0}
\qfigdef{p1_phy_s2_q07_1}{Score-II - Test 40 (16 Mar 2025)-img/p1_phy_s2_q07_1.png}{81.2}{0}
\qfigdef{p1_phy_s3_q11_1}{Score-II - Test 40 (16 Mar 2025)-img/p1_phy_s3_q11_1.png}{97.1}{0}
\qfigdef{p1_phy_s4_q02_1}{Score-II - Test 40 (16 Mar 2025)-img/p1_phy_s4_q02_1.png}{156.2}{0}
\qfigdef{p1_phy_s4_q06_1}{Score-II - Test 40 (16 Mar 2025)-img/p1_phy_s4_q06_1.png}{88.8}{0}
\qfigdef{p1_chem_s1_q01_1}{Score-II - Test 40 (16 Mar 2025)-img/p1_chem_s1_q01_1.png}{210.1}{0}
\qfigdef{p1_chem_s1_q04_1}{Score-II - Test 40 (16 Mar 2025)-img/p1_chem_s1_q04_1.png}{134.5}{0}
\qfigdef{p1_chem_s2_q06_1}{Score-II - Test 40 (16 Mar 2025)-img/p1_chem_s2_q06_1.png}{169.6}{0}
\qfigdef{p1_chem_s2_q07_1}{Score-II - Test 40 (16 Mar 2025)-img/p1_chem_s2_q07_1.png}{204.2}{0}
\qfigdef{p1_chem_s2_q07_2}{Score-II - Test 40 (16 Mar 2025)-img/p1_chem_s2_q07_2.png}{122.3}{0}
\qfigdef{p1_chem_s2_q07_A}{Score-II - Test 40 (16 Mar 2025)-img/p1_chem_s2_q07_A.png}{135.9}{0}
\qfigdef{p1_chem_s2_q07_B}{Score-II - Test 40 (16 Mar 2025)-img/p1_chem_s2_q07_B.png}{133.0}{0}
\qfigdef{p1_chem_s2_q07_C}{Score-II - Test 40 (16 Mar 2025)-img/p1_chem_s2_q07_C.png}{127.0}{0}
\qfigdef{p1_chem_s2_q07_D}{Score-II - Test 40 (16 Mar 2025)-img/p1_chem_s2_q07_D.png}{132.7}{0}
\qfigdef{p1_chem_s3_q10_4}{Score-II - Test 40 (16 Mar 2025)-img/p1_chem_s3_q10_4.png}{119.9}{0}
\qfigdef{p1_chem_s4_q03_1}{Score-II - Test 40 (16 Mar 2025)-img/p1_chem_s4_q03_1.png}{138.7}{0}
\qfigdef{p1_chem_s4_q04_1}{Score-II - Test 40 (16 Mar 2025)-img/p1_chem_s4_q04_1.png}{157.1}{0}
\qfigdef{p1_chem_s4_q05_1}{Score-II - Test 40 (16 Mar 2025)-img/p1_chem_s4_q05_1.png}{134.1}{0}
\qfigdef{p1_chem_s4_q06_1}{Score-II - Test 40 (16 Mar 2025)-img/p1_chem_s4_q06_1.png}{123.5}{0}
\qfigdef{p2_phy_s1_q01_1}{Score-II - Test 40 (16 Mar 2025)-img/p2_phy_s1_q01_1.png}{171.3}{0}
\qfigdef{p2_phy_s2_q05_1}{Score-II - Test 40 (16 Mar 2025)-img/p2_phy_s2_q05_1.png}{117.2}{0}
\qfigdef{p2_phy_s2_q06_1}{Score-II - Test 40 (16 Mar 2025)-img/p2_phy_s2_q06_1.png}{144.9}{0}
\qfigdef{p2_phy_s3_para1_2_1}{Score-II - Test 40 (16 Mar 2025)-img/p2_phy_s3_para1_2_1.png}{84.0}{0}
\qfigdef{p2_phy_s3_para3_4_1}{Score-II - Test 40 (16 Mar 2025)-img/p2_phy_s3_para3_4_1.png}{98.6}{0}
\qfigdef{p2_phy_s4_q05_1}{Score-II - Test 40 (16 Mar 2025)-img/p2_phy_s4_q05_1.png}{145.4}{0}
\qfigdef{p2_phy_s4_q06_1}{Score-II - Test 40 (16 Mar 2025)-img/p2_phy_s4_q06_1.png}{107.3}{0}
\qfigdef{p2_phy_s4_q08_1}{Score-II - Test 40 (16 Mar 2025)-img/p2_phy_s4_q08_1.png}{125.8}{0}
\qfigdef{p2_phy_s4_q10_1}{Score-II - Test 40 (16 Mar 2025)-img/p2_phy_s4_q10_1.png}{233.8}{0}
\qfigdef{p2_chem_s1_q01_A}{Score-II - Test 40 (16 Mar 2025)-img/p2_chem_s1_q01_A.png}{102.6}{0}
\qfigdef{p2_chem_s1_q01_B}{Score-II - Test 40 (16 Mar 2025)-img/p2_chem_s1_q01_B.png}{103.7}{0}
\qfigdef{p2_chem_s1_q01_C}{Score-II - Test 40 (16 Mar 2025)-img/p2_chem_s1_q01_C.png}{102.5}{0}
\qfigdef{p2_chem_s1_q01_D}{Score-II - Test 40 (16 Mar 2025)-img/p2_chem_s1_q01_D.png}{102.9}{0}
\qfigdef{p2_chem_s1_q03_1}{Score-II - Test 40 (16 Mar 2025)-img/p2_chem_s1_q03_1.png}{132.7}{0}
\qfigdef{p2_chem_s1_q04_A}{Score-II - Test 40 (16 Mar 2025)-img/p2_chem_s1_q04_A.png}{151.2}{0}
\qfigdef{p2_chem_s1_q04_B}{Score-II - Test 40 (16 Mar 2025)-img/p2_chem_s1_q04_B.png}{151.6}{0}
\qfigdef{p2_chem_s1_q04_C}{Score-II - Test 40 (16 Mar 2025)-img/p2_chem_s1_q04_C.png}{161.9}{0}
\qfigdef{p2_chem_s1_q04_D}{Score-II - Test 40 (16 Mar 2025)-img/p2_chem_s1_q04_D.png}{161.5}{0}
\qfigdef{p2_chem_s2_q07_1}{Score-II - Test 40 (16 Mar 2025)-img/p2_chem_s2_q07_1.png}{184.4}{0}
\qfigdef{p2_chem_s3_para1_2_1}{Score-II - Test 40 (16 Mar 2025)-img/p2_chem_s3_para1_2_1.png}{121.8}{0}
\qfigdef{p2_math_s3_para3_4_1}{Score-II - Test 40 (16 Mar 2025)-img/p2_math_s3_para3_4_1.png}{76.4}{0}
\raggedcolumns
\makeatletter\@beginparpenalty=10000\makeatother
\begin{document}
\begin{center}{\Large\bfseries JEE (Main+Advanced) Score-II --- Test 40 (16 Mar 2025)}\end{center}
\begin{multicols}{2}

\papersec{1}
\partsec{1}{PHYSICS}
\sect{Section-I(i)}

\qn{1} A bar of mas $M=1.00$ kg and length $L=0.20$ m is lying on a horizontal frictionless surface. A small mass $m=0.10$ kg is moving on the same horizontal surface with 5.00 m s$^{-1}$ speed on a path perpendicular to the bar. It hits the bar at a distance $L/4$ from the left end and returns back on the same path with speed $v$. After this elastic collision, the bar rotates with an angular velocity $\omega$. Which of the following statement is correct?
\begin{qoptions}
\opt{A}{$\omega=\dfrac{600}{47}$ rad s$^{-1}$ and $v=\dfrac{165}{47}$ m s$^{-1}$}
\opt{B}{$\omega=\dfrac{300}{47}$ rad s$^{-1}$ and $v=\dfrac{165}{47}$ m s$^{-1}$}
\opt{C}{$\omega=\dfrac{200}{23}$ rad s$^{-1}$ and $v=\dfrac{51}{23}$ m s$^{-1}$}
\opt{D}{$\omega=\dfrac{800}{23}$ rad s$^{-1}$ and $v=\dfrac{51}{23}$ m s$^{-1}$}
\end{qoptions}

\qn{2} A container has a base of 50 cm $\times$ 5 cm and height 50 cm, as shown in the figure. It has two parallel electrically conducting walls each of area 50 cm $\times$ 50 cm. The remaining walls of the container are thin and non-conducting. The container is being filled with a liquid of dielectric constant 3 at a uniform rate of 250 cm$^3$ s$^{-1}$. What is the value of the capacitance of the container after 20 seconds ?\\
\relax[Given: Permittivity of free space $\epsilon_0=9\times10^{-12}\ \text{C}^2\,\text{N}^{-1}\,\text{m}^{-2}$, the effects of the non-conducting walls on the capacitance are negligible]
\qfig{p1_phy_s1_q02_1}
\begin{qoptions}
\opt{A}{27 pF}
\opt{B}{63 pF}
\opt{C}{81 pF}
\opt{D}{135 pF}
\end{qoptions}

\qn{3} One mole of an ideal gas expands adiabatically from an initial state ($P_0$, $V_0$) to final state ($P_A$, $5V_0$). Another mole of the same gas expands isothermally from a same initial state ($P_0$, $V_0$) to the vol. ($P_B$, $5V_0$). The ratio of the specific heats at constant pressure and constant volume of this ideal gas is $\gamma$. What is the ratio $\dfrac{P_A}{P_B}$ ?
\begin{qoptions}
\opt{A}{$5^{\gamma-1}$}
\opt{B}{$5^{1-\gamma}$}
\opt{C}{$5^\gamma$}
\opt{D}{$5^{1+\gamma}$}
\end{qoptions}

\qn{4} Two satellites P and Q are moving in different circular orbits around the Earth (radius $R$). The heights of P and Q from the Earth surface are $h_P$ and $h_Q$, respectively, where $h_P=R/3$. The accelerations of P and Q due to Earth's gravity are $g_P$ and $g_Q$, respectively. If $g_P/g_Q=36/25$, what is the velocity of Q?\\
(M = Mass of earth, G = Universal gravitational constant)
\begin{qoptions}
\opt{A}{$\sqrt{\dfrac{5GM}{3R}}$}
\opt{B}{$\sqrt{\dfrac{5GM}{8R}}$}
\opt{C}{$\sqrt{\dfrac{5GM}{6R}}$}
\opt{D}{$\sqrt{\dfrac{6GM}{7R}}$}
\end{qoptions}

\sect{Section-I(ii)}

\qn{5} A slide with a frictionless curved surface, which becomes horizontal at its lower end,, is fixed on the terrace of a building of height 2h from the ground, as shown in the figure. A spherical ball of mass m is released on the slide from rest at a height h from the top of the terrace. The ball leaves the slide with a velocity $\vec u_0=u_0\hat x$ and falls on the ground at a distance d from the building making an angle $\theta$ with the horizontal. It bounces off with a velocity $\vec v$ and reaches a maximum height $h_1$. The acceleration due to gravity is g and the collision is elastic. Which of the following statement(s) is(are) correct ? [Take h = 20 m, g = 10 m/s$^2$]. The ground is frictionless.
\qfig{p1_phy_s2_q05_1}
\begin{qoptions}
\opt{A}{$\vec u_0=20\hat i$ m/s}
\opt{B}{$\vec v=\left(20\hat i+20\sqrt2\hat k\right)$ m/s}
\opt{C}{$\theta=60^\circ$}
\opt{D}{$\dfrac d{h_1}=\sqrt2$}
\end{qoptions}

\qn{6} A plane polarized blue light ray is incident on a prism such that there is no reflection from the surface of the prism. The angle of deviation of the emergent ray is $\delta=60^\circ$ (see Figure-1). The angle of minimum deviation for red light from the same prism is $\delta_{min}=30^\circ$ (see Figure-2). The refractive index of the prism material for blue light is $\sqrt3$
\qfig{p1_phy_s2_q06_2}
\qfig{p1_phy_s2_q06_1}
\begin{qoptions}
\opt{A}{The blue light has its electric field in the plane of incidence}
\opt{B}{The angle of the prism is $60^\circ$.}
\opt{C}{The refractive index of the material of the prism for red light is $\sqrt2$.}
\opt{D}{The angle of refraction for blue light in air at the exit plane of the prism is $45^\circ$.}
\end{qoptions}

\qn{7} In a circuit shown in the figure, the capacitor $C$ is initially uncharged and the key $K$ is open. In this condition, a current of 3A flows through 1 $\Omega$ resistor. The key is closed at time $t=t_0$. Which of the following statement(s) is(are) correct? [Given: $e^{-1}\approx0.35$]
\qfig{p1_phy_s2_q07_1}
\begin{qoptions}
\opt{A}{The value of the resistance $R$ is 1$\Omega$.}
\opt{B}{Long after the switch is closed, value of current $I_1$ is 4 A.}
\opt{C}{At $t=t_0+8\ \mu$s, the current in the capacitor branch is approxmately 0.7 A.}
\opt{D}{For $t\to\infty$, the charge on the capacitor is 16$\mu C$.}
\end{qoptions}

\sect{Section-I(iii)}

\qn{8} List-I shows different radioactive decay processes and List-II provides possible emitted particles. Match each entry in List-I with an appropriate entry from List-II, and choose the correct option.\par\nobreak
\begin{center}\small\setlength{\tabcolsep}{3pt}
\begin{tabular}{|c|p{2.6cm}|c|p{3.6cm}|}
\hline
 & \multicolumn{1}{c|}{\textbf{List$-$I}} & & \multicolumn{1}{c|}{\textbf{List$-$II}} \\
\hline
(P) & ${}^{238}_{92}\text{U}\to{}^{234}_{91}\text{Pa}$ & (1) & three $\beta^-$ particles and one $\alpha$ particle \\
\hline
(Q) & ${}^{214}_{82}\text{Pb}\to{}^{210}_{82}\text{Pb}$ & (2) & two $\beta^-$ particles and one $\alpha$ particle \\
\hline
(R) & ${}^{210}_{81}\text{Tl}\to{}^{206}_{82}\text{Pb}$ & (3) & one $\alpha$ particle and one $\beta^+$ particle \\
\hline
(S) & ${}^{228}_{91}\text{Pa}\to{}^{224}_{88}\text{Ra}$ & (4) & one $\alpha$ particle and two $\beta^+$ particles \\
\hline
 & & (5) & one $\alpha$ particle and one $\beta^-$ particle \\
\hline
\end{tabular}
\end{center}
\begin{qoptions}
\opt{A}{P $\to$ 5; Q $\to$ 2; R $\to$ 1; S $\to$ 3}
\opt{B}{P $\to$ 4; Q $\to$ 3; R $\to$ 4; S $\to$ 2}
\opt{C}{P $\to$ 5; Q $\to$ 3; R $\to$ 1; S $\to$ 4}
\opt{D}{P $\to$ 5; Q $\to$ 2; R $\to$ 4; S $\to$ 1}
\end{qoptions}

\qn{9} Match the temperature of a black body given in List-I with an appropriate statement in List-II, and choose the correct option. [Given: Wien's constant as $2.9\times10^{-3}$ m-K and $\dfrac{hc}e=1.24\times10^{-6}$ V-m]\par\nobreak
\begin{center}\small\setlength{\tabcolsep}{3pt}
\begin{tabular}{|c|p{1.1cm}|c|p{4.6cm}|}
\hline
 & \multicolumn{1}{c|}{\textbf{List-I}} & & \multicolumn{1}{c|}{\textbf{List-II}} \\
\hline
(P) & 2000 K & (1) & The radiation at peak emission wavelength will result in the widest central maximum of a single slit diffraction. \\
\hline
(Q) & 3000 K & (2) & The power emitted per unit area is 1/16 of that emitted by a blackbody at temperature 6000 K. \\
\hline
(R) & 5000 K & (3) & The radiation at peak wavelength is visible to human eye. \\
\hline
(S) & 10000 K & (4) & The radiation at peak emission wavelength can be used to image human bones. \\
\hline
 & & (5) & The radiation at peak wavelength can lead to emission of photoelectrons from a metal of work function 4 eV \\
\hline
\end{tabular}
\end{center}
\begin{qoptions}
\opt{A}{P $\to$ 4; Q $\to$ 2; R $\to$ 1; S $\to$ 5}
\opt{B}{P $\to$ 3; Q $\to$ 2; R $\to$ 4; S $\to$ 1}
\opt{C}{P $\to$ 1; Q $\to$ 2; R $\to$ 3; S $\to$ 5}
\opt{D}{P $\to$ 1; Q $\to$ 2; R $\to$ 5; S $\to$ 3}
\end{qoptions}

\qn{10} A series LCR circuit is connected to a 45 sin ($\omega t$) Volt source. The resonant angular frequency of the circuit is $10^5$ rad s$^{-1}$ and current amplitude at resonance is $I_0$. When the angular frequency of the source is $\omega=8\times10^4$ rad s$^{-1}$, the current amplitude in the circuit is 0.05 $I_0$. If $L$ = 50 mH, match each entry in List-I with an appropriate value from List-II and choose the correct option.\par\nobreak
\begin{center}\small\setlength{\tabcolsep}{3pt}
\begin{tabular}{|c|p{3.9cm}|c|p{0.8cm}|}
\hline
 & \multicolumn{1}{c|}{\textbf{List$-$I}} & & \multicolumn{1}{c|}{\textbf{\footnotesize List$-$II}} \\
\hline
(P) & $I_0$ in mA & (1) & 2250 \\
\hline
(Q) & The quality factor of the circuit & (2) & 500 \\
\hline
(R) & The bandwidth of the circuit in rad s$^{-1}$ & (3) & 18 \\
\hline
(S) & The peak power dissipated at resonance in Watt & (4) & 400 \\
\hline
 & & (5) & 44.4 \\
\hline
\end{tabular}
\end{center}
\begin{qoptions}
\opt{A}{P $\to$ 3; Q $\to$ 1; R $\to$ 2; S $\to$ 4}
\opt{B}{P $\to$ 3; Q $\to$ 1; R $\to$ 4; S $\to$ 2}
\opt{C}{P $\to$ 4; Q $\to$ 5; R $\to$ 3; S $\to$ 1}
\opt{D}{P $\to$ 4; Q $\to$ 5; R $\to$ 1; S $\to$ 3}
\end{qoptions}

\qn{11} A thin conducting rod MN of mass 20 gm, length 25 cm and resistance 10 $\Omega$ is held on frictionless, long, perfectly conducting vertical rails as shown in the figure. There is a uniform magnetic field $B_0$ = 4 T directed perpendicular to the plane of the rod-rail arrangement. The rod is released from rest at time t = 0 and it moves down along the rails. Assume air drag is negligible. Match each quantity in List-I with an appropriate value from List-II, and choose the correct option. [Given : The acceleration due to gravity $g=10\ \text{ms}^{-2}$ and $e^{-1}=0.4$]
\qfig{p1_phy_s3_q11_1}\par\nobreak
\begin{center}\small\setlength{\tabcolsep}{3pt}
\begin{tabular}{|c|p{3.6cm}|c|p{0.8cm}|}
\hline
 & \multicolumn{1}{c|}{\textbf{List$-$I}} & & \multicolumn{1}{c|}{\textbf{\footnotesize List$-$II}} \\
\hline
(P) & At t = 0.2 s, the magnitude of the induced emf in Volt & (1) & 1.20 \\
\hline
(Q) & At t = 0.2 s, the magnitude of the magnetic force in Newton & (2) & 0.07 \\
\hline
(R) & At t = 0.2 s, the power dissipated as heat in Watt & (3) & 2.00 \\
\hline
(S) & The magnitude of terminal velocity of the rod in m s$^{-1}$ & (4) & 0.12 \\
\hline
 & & (5) & 0.14 \\
\hline
\end{tabular}
\end{center}
\begin{qoptions}
\opt{A}{P $\to$ 5; Q $\to$ 2; R $\to$ 3; S $\to$ 1}
\opt{B}{P $\to$ 1; Q $\to$ 4; R $\to$ 5; S $\to$ 3}
\opt{C}{P $\to$ 4; Q $\to$ 3; R $\to$ 1; S $\to$ 2}
\opt{D}{P $\to$ 4; Q $\to$ 1; R $\to$ 5; S $\to$ 3}
\end{qoptions}

\sect{Section-II}

\qn{1} A Hydrogen-like atom has atomic number $Z$. Photons emitted in the electronic transitions from level $n=5$ to level $n=3$ in these atoms are used to perform photoelectric effect experiment on a target metal. The maximum kinetic energy of the photoelectrons generated is 1.86 eV. If the photoelectric threshold wavelength for the target metal is 620 nm, the value of $Z$ is \underline{\hspace{1cm}}. [Given: $hc=1240$ eV-nm and $Rhc=13.6$ eV, where $R$ is the Rydberg constant, $h$ is the Planck's constant and $c$ is the speed of light in vacuum]

\qn{2} An optical arrangement consists of two concave mirrors $M_1$ and $M_2$, and a convex lens L with a common principal axis, as shown in the figure. The focal length of L is 10 cm. The radii of curvature of $M_1$ and $M_2$ are 20 cm and 24 cm, respectively. The distance between L and $M_2$ is 40 cm. A point object S is placed at the mid-point between L and $M_2$ on the axis. When the distance between L and $M_1$ is x cm, one of the images coincides with S. The value of $x$ is \underline{\hspace{1cm}}. (Consider at least two reflections)
\qfig{p1_phy_s4_q02_1}

\qn{3} In an experiment for determination of the focal length of a thin convex lens, the distance of the object from the lens is $10\pm0.1$ cm and the distance of its real image from the lens is $15\pm0.1$ cm. The error in the determination of focal length of the lens is \underline{\hspace{1cm}} $\mu$m.

\qn{4} A closed container contains a homogeneous mixture of one mole of an ideal monatomic gas ($\gamma=5/3$) and one mole of an ideal diatomic gas ($\gamma=7/5$). Here, $\gamma$ is the ratio of the specific heats at constant pressure and constant volume of an ideal gas. The gas mixture does a work of 66 Joule when heated at constant pressure. The change in its internal energy is \underline{\hspace{1cm}} Joule.

\qn{5} A person of height 1.6 m is walking away from a lamp post of height 4 m along a straight path on the flat ground. The lamp post and the person are always perpendicular to the ground. If the speed of the person is 60 cm s$^{-1}$, the speed of the tip of the person's shadow on the ground with respect to the ground is \underline{\hspace{1cm}} cm s$^{-1}$.

\qn{6} Two point-like objects of masses 20 gm and 30 gm are fixed at the two ends of a rigid massless rod of length 10 cm. This system is suspended vertically from a rigid ceiling using a thin wire attached to its center of mass, as shown in the figure. The resulting torsional pendulum undergoes small oscillations of amplitude 0.01 rad. The torsional constant of the wire is $1.2\times10^{-8}$ N m rad$^{-1}$. The maximum speed of 20 gm mass is \underline{\hspace{1cm}} $\mu$ m/s.
\qfig{p1_phy_s4_q06_1}

\partsec{2}{CHEMISTRY}
\sect{Section-I(i)}

\qn{1} In the scheme given below, \textbf{X, Y} and \textbf{P} respectively, are
\qfig{p1_chem_s1_q01_1}
\begin{qoptions}
\opt{A}{$\text{CrO}_4^{2-}$, Br$_2$ and MnO}
\opt{B}{$\text{MnO}_4^{2-}$, Cl$_2$ and Mn(OH)$_2$}
\opt{C}{$\text{MnO}_4^-$, Cl$_2$ and Mn(OH)$_2$}
\opt{D}{MnSO$_4$, HOCl and MnO}
\end{qoptions}

\qn{2} Plotting $1/\Lambda_m$ (Y-axis) against $c\Lambda_m$ (X-axis) for aqueous solutions of a monobasic weak acid (HX) resulted in a straight line with y-axis intercept of P and slope of S. The $K_a$ is\\
$\left[\Lambda_m=\text{molar conductivity}\right.$\\
$\Lambda_m^o=$ limiting molar conductivity\\
c = molar concentration\\
$\left.K_a=\text{dissociation constant of HX}\right]$
\begin{qoptions}
\opt{A}{$\dfrac PS$}
\opt{B}{$\dfrac SP$}
\opt{C}{$\dfrac{P^2}S$}
\opt{D}{$\dfrac P{S^2}$}
\end{qoptions}

\qn{3} On decreasing the $p$H from 7 to 2, the solubility of a sparingly soluble salt (MX) of a weak acid (HX) increased from $10^{-4}$ mol L$^{-1}$ to `x' mol L$^{-1}$. Then `x' is ($K_a=10^{-4}$) :
\begin{qoptions}
\opt{A}{$10^{-3}$}
\opt{B}{$10^{-4}$}
\opt{C}{$10^{-5}$}
\opt{D}{$10^{-2}$}
\end{qoptions}

\qn{4} In the given reaction scheme, \textbf{P} is a phenyl alkyl ether; \textbf{Q} is an aromatic compound; \textbf{R} and \textbf{S} are the major products.
\qfig{p1_chem_s1_q04_1}
The correct statement about \textbf{S} is
\begin{qoptions}
\opt{A}{It primarily inhibits noradrenaline degrading enzymes.}
\opt{B}{It inhibits the synthesis of prostaglandin.}
\opt{C}{It is a narcotic drug.}
\opt{D}{It is \textit{ortho}-acetylbenzoic acid.}
\end{qoptions}

\sect{Section-I(ii)}

\qn{5} The correct statement(s) related to processes involved in the extraction of metals is(are)
\begin{qoptions}
\opt{A}{Zone refining of metals is based on the principle of greater mobility of pure metal than that of impurity}
\opt{B}{Calcination of calamine produces sphalerite}
\opt{C}{Copper pyrite is heated with silica in a reverberatory furnace to remove iron}
\opt{D}{When Malachite is heated in a reverberatory furnace, cupric oxide is formed.}
\end{qoptions}

\qn{6} In the following reactions, \textbf{P}, \textbf{Q}, \textbf{R}, and \textbf{S} are the major products.
\qfig{p1_chem_s2_q06_1}
The correct statement(s) about \textbf{P}, \textbf{Q}, \textbf{R}, and \textbf{S} is(are)
\begin{qoptions}
\opt{A}{Both \textbf{P} and \textbf{Q} have asymmetric carbon(s).}
\opt{B}{Both \textbf{Q} and \textbf{R} have asymmetric carbon(s).}
\opt{C}{Both \textbf{P} and \textbf{R} have asymmetric carbon(s).}
\opt{D}{\textbf{P} has asymmetric carbon(s), \textbf{S} does \textbf{not} have any asymmetric carbon.}
\end{qoptions}

\qn{7} Consider the following reaction scheme and choose the correct option(s) for the major products \textbf{Q}, \textbf{R} and \textbf{S}.
\qfig{p1_chem_s2_q07_1}
\qfig{p1_chem_s2_q07_2}
\begin{qoptions}
\opt{A}{\qfigi{p1_chem_s2_q07_A}}
\opt{B}{\qfigi{p1_chem_s2_q07_B}}
\opt{C}{\qfigi{p1_chem_s2_q07_C}}
\opt{D}{\qfigi{p1_chem_s2_q07_D}}
\end{qoptions}

\sect{Section-I(iii)}

\qn{8} Match the reactions (in the given stoichiometry of the reactants) in List-I with one of their products given in List-II and choose the correct option.\par\nobreak
\begin{center}\small\setlength{\tabcolsep}{3pt}
\begin{tabular}{|c|p{3.4cm}|c|p{1.6cm}|}
\hline
 & \multicolumn{1}{c|}{\textbf{List-I}} & & \multicolumn{1}{c|}{\textbf{List-II}} \\
\hline
(P) & PCl$_3$ + H$_3$PO$_3$ $\rightarrow$ & (1) & H$_3$PO$_2$ \\
\hline
(Q) & Red P$_4$ + alkali $\rightarrow$ & (2) & H$_4$P$_2$O$_5$ \\
\hline
(R) & C$_2$H$_5$OH + PCl$_3$ $\rightarrow$ & (3) & H$_4$P$_2$O$_6$ \\
\hline
(S) & H$_3$PO$_2$ + 2H$_2$O + 4AgNO$_3$ $\rightarrow$ & (4) & H$_3$PO$_3$ \\
\hline
 & & (5) & H$_3$PO$_4$ \\
\hline
\end{tabular}
\end{center}
\begin{qoptions}
\opt{A}{P $\to$ 2; Q $\to$ 3; R $\to$ 1; S $\to$ 5}
\opt{B}{P $\to$ 3; Q $\to$ 5; R $\to$ 4; S $\to$ 2}
\opt{C}{P $\to$ 5; Q $\to$ 2; R $\to$ 1; S $\to$ 3}
\opt{D}{P $\to$ 2; Q $\to$ 3; R $\to$ 4; S $\to$ 5}
\end{qoptions}

\qn{9} Match the electronic configurations in List-I with appropriate metal complex ions in List-II and choose the correct option.\\
\relax[Atomic Number: Fe = 26, Mn = 25, Co = 27]\par\nobreak
\begin{center}\small\setlength{\tabcolsep}{3pt}
\begin{tabular}{|c|p{1.5cm}|c|p{2.4cm}|}
\hline
 & \multicolumn{1}{c|}{\textbf{List-I}} & & \multicolumn{1}{c|}{\textbf{List-II}} \\
\hline
(P) & $t_{2g}^6e_g^0$ & (1) & $\left[\text{NiF}_6\right]^{-2}$ \\
\hline
(Q) & $t_{2g}^6e_g^3$ & (2) & $\left[\text{CoF}_6\right]^{3-}$ \\
\hline
(R) & $e^2t_2^3$ & (3) & $\left[\text{Cu(H}_2\text{O)}_6\right]^{2+}$ \\
\hline
(S) & $t_{2g}^4e_g^2$ & (4) & $\left[\text{FeCl}_4\right]^{-}$ \\
\hline
 & & (5) & $\left[\text{CoCl}_4\right]^{2-}$ \\
\hline
\end{tabular}
\end{center}
\begin{qoptions}
\opt{A}{P $\to$ 1; Q $\to$ 4; R $\to$ 2; S $\to$ 3}
\opt{B}{P $\to$ 1; Q $\to$ 2; R $\to$ 4; S $\to$ 5}
\opt{C}{P $\to$ 3; Q $\to$ 2; R $\to$ 5; S $\to$ 1}
\opt{D}{P $\to$ 1; Q $\to$ 3; R $\to$ 4; S $\to$ 2}
\end{qoptions}

\qn{10} Match the reactions in List-I with the features of their products in List-II and choose the correct option.\par\nobreak
\begin{center}\small\setlength{\tabcolsep}{3pt}
\begin{tabular}{|c|p{3.6cm}|c|p{2.0cm}|}
\hline
\multicolumn{2}{|c|}{\textbf{List-I}} & \multicolumn{2}{c|}{\textbf{List-II}} \\
\hline
(P) & (-)-1-Bromo-2-ethylpentane $\xrightarrow[\text{S}_\text{N}\text{2 reaction}]{\text{aq. NaOH}}$ (single enantiomer) & (1) & Inversion of configuration \\
\hline
(Q) & (-)-2-Bromopentane $\xrightarrow[\text{S}_\text{N}\text{2 reaction}]{\text{aq. NaOH}}$ (single enantiomer) & (2) & Retention of configuration \\
\hline
(R) & (-)-3-Bromo-3-methylhexane $\xrightarrow[\text{S}_\text{N}\text{1 reaction}]{\text{aq. NaOH}}$ (single enantiomer) & (3) & Mixture of enantiomers \\
\hline
(S) & \qfigi{p1_chem_s3_q10_4} & (4) & Mixture of structural isomers \\
\hline
 & & (5) & Mixture of diastereomers \\
\hline
\end{tabular}
\end{center}
\begin{qoptions}
\opt{A}{P $\to$ 1; Q $\to$ 2; R $\to$ 5; S $\to$ 3}
\opt{B}{P $\to$ 2; Q $\to$ 1; R $\to$ 3; S $\to$ 5}
\opt{C}{P $\to$ 1; Q $\to$ 2; R $\to$ 5; S $\to$ 4}
\opt{D}{P $\to$ 2; Q $\to$ 4; R $\to$ 3; S $\to$ 5}
\end{qoptions}

\qn{11} The major products obtained from the reactions in List-II are the reactants for the named reactions mentioned in List-I. Match List-I with List-II and choose the correct option.\par\nobreak
\begin{center}\small\setlength{\tabcolsep}{3pt}
\begin{tabular}{|c|p{2.3cm}|c|p{3.0cm}|}
\hline
 & \multicolumn{1}{c|}{\textbf{List-I}} & & \multicolumn{1}{c|}{\textbf{List-II}} \\
\hline
(P) & Etard reaction & (1) & Acetophenone $\xrightarrow{\text{Zn--Hg, HCl}}$ \\
\hline
(Q) & Gattermann reaction & (2) & Toluene $\xrightarrow[\text{(ii) SOCl}_2]{\text{(i) KMnO}_4,\text{KOH, }\Delta}$ \\
\hline
(R) & Gattermann-Koch reaction & (3) & Benzene $\xrightarrow[\text{anhyd. AlCl}_3]{\text{CH}_3\text{Cl}}$ \\
\hline
(S) & Rosenmund reduction & (4) & Aniline $\xrightarrow[\text{273--278 K}]{\text{NaNO}_2/\text{HCl}}$ \\
\hline
 & & (5) & Phenol $\xrightarrow{\text{Zn, }\Delta}$ \\
\hline
\end{tabular}
\end{center}
\begin{qoptions}
\opt{A}{P $\to$ 2; Q $\to$ 4; R $\to$ 1; S $\to$ 3}
\opt{B}{P $\to$ 1; Q $\to$ 3; R $\to$ 5; S $\to$ 2}
\opt{C}{P $\to$ 3; Q $\to$ 2; R $\to$ 1; S $\to$ 4}
\opt{D}{P $\to$ 3; Q $\to$ 4; R $\to$ 5; S $\to$ 2}
\end{qoptions}

\sect{Section-II}

\qn{1} The stoichiometric reaction of 516 g of dimethyldichlorosilane with water results in a tetrameric cyclic product \textbf{X} in 50\% yield. The weight (in g) of \textbf{X} obtained is\underline{\hspace{1cm}}.\\
\relax[Use, molar mass (g mol$^{-1}$): H = 1, C = 12, O = 16, Si = 28, Cl = 35.5]

\qn{2} A gas has molar volume of 0.4 dm$^3$ mol$^{-1}$ at a temperature of 800 K and pressure \textbf{x} atm. If it shows ideal gas behaviour at the same temperature and pressure, the molar volume is 0.8 dm$^3$ mol$^{-1}$. The value of \textbf{x} is \underline{\hspace{1cm}}.\\
\relax[Use: Gas constant, R = $8\times10^{-2}$ L atm K$^{-1}$ mol$^{-1}$]

\qn{3} The plot of log $k_f$ versus $\dfrac1T$ for a reversible reaction A (g) $\rightleftharpoons$ P (g) is shown.
\qfig{p1_chem_s4_q03_1}
Pre-exponential factors for the forward and backward reactions are $10^{15}$ s$^{-1}$ and $10^{11}$ s$^{-1}$, respectively. If the value of log $K_b$ for the reaction at 250 K is $-5$, the value of log K at 500 K is \underline{\hspace{1cm}}.\\
\relax[K = equilibrium constant of the reaction\\
$k_f$ = rate constant of forward reaction\\
$k_b$ = rate constant of backward reaction]

\qn{4} One mole of an ideal monatomic gas undergoes two reversible processes (A $\to$ B and B $\to$ C) as shown in the given figure :
\qfig{p1_chem_s4_q04_1}
A $\to$ B is an adiabatic process. If the total entropy change in the entire process (A $\to$ B and B $\to$ C) is R ln $\sqrt{10}$, the value of log $V_3$ is \underline{\hspace{1cm}}.\\
\relax[Use, molar heat capacity of the gas at constant pressure, $C_{p,m}=\dfrac52$ R]

\qn{5} In a one-litre flask, 6 moles of A undergoes the reaction A(g) $\rightleftharpoons$ P(g). The progress of product formation at two temperatures (in Kelvin), $T_1$ and $T_2$, is shown in the figure:
\qfig{p1_chem_s4_q05_1}
If $T_1=2T_2$ and $\Delta H^0=RT_2\ln x$, then the value of x is \underline{\hspace{1cm}}.

\qn{6} The total number of sp$^2$ hybridised carbon atoms in the major product \textbf{P} (a non-heterocyclic compound) of the following reaction is \underline{\hspace{1cm}}.
\qfig{p1_chem_s4_q06_1}

\partsec{3}{MATHEMATICS}
\sect{Section-I(i)}

\qn{1} Let $f:(0,1)\to\mathbb R$ be the function defined as $f(x)=\sqrt n$ if $x\in\left[\dfrac1{n+1},\dfrac1n\right)$ where $n\in\mathbb N$.\\
Let $g:(0,1)\to\mathbb R$ be a function such that $\displaystyle\int_{x^2}^x\sqrt{\dfrac{1-t}t}\,dt<g(x)<2\sqrt x$ for all $x\in(0,1)$. Then $\displaystyle\lim_{x\to0}f(x)g(x)$
\begin{qoptions}
\opt{A}{is equal to 2}
\opt{B}{is equal to 1}
\opt{C}{does \textbf{NOT} exist}
\opt{D}{is equal to 3}
\end{qoptions}

\qn{2} Let Q be the cube with the set of vertices $\left\{(x_1,x_2,x_3)\in\mathbb R^3:x_1,x_2,x_3\in\{-1,1\}\right\}$. Let F be the set of all twelve lines containing the diagonals of the six faces of the cube Q. Let S be the set of all four lines containing the main diagonals of the cube Q; for instance, the line passing through the vertices $(-1,-1,-1)$ and $(1,1,1)$ is in S. For lines $\ell_1$ and $\ell_2$, let d($\ell_1$, $\ell_2$) denote the shortest distance between them. Then the maximum value of d($\ell_1$, $\ell_2$), as $\ell_1$ varies over F and $\ell_2$ varies over S, is
\begin{qoptions}
\opt{A}{$\dfrac1{\sqrt6}$}
\opt{B}{$\dfrac{\sqrt2}{\sqrt3}$}
\opt{C}{$\dfrac1{\sqrt3}$}
\opt{D}{$\dfrac1{\sqrt{12}}$}
\end{qoptions}

\qn{3} Let\\ $X=\left\{(x,y)\in\mathbb Z\times\mathbb Z:\dfrac{(x+1)^2}8+\dfrac{(y+3)^2}{20}<1\right.$\\ $\left.\text{and }(y+3)^2<5(x+1)\right\}$. Three distinct points P, Q and R are randomly chosen from X. Then the probability that P, Q and R form a triangle whose area is a positive integer, is
\begin{qoptions}
\opt{A}{$\dfrac{83}{220}$}
\opt{B}{$\dfrac{71}{220}$}
\opt{C}{$\dfrac{73}{220}$}
\opt{D}{$\dfrac{79}{220}$}
\end{qoptions}

\qn{4} Let P be a point on the parabola $y^2=4ax$, where $a>0$. The normal to the parabola at P meets the x-axis at a point Q. The area of the triangle PFQ, where F is the focus of the parabola, is 90. If the slope m of the normal and a are both positive integers, then the pair (a, m) is
\begin{qoptions}
\opt{A}{(2, 3)}
\opt{B}{(3, 2)}
\opt{C}{(2, 4)}
\opt{D}{(4, 2)}
\end{qoptions}

\sect{Section-I(ii)}

\qn{5} Let $S=(-2,-1)\cup(0,1)\cup(2,3)\cup(4,5)$ and $T=\{0,1,2,3,4\}$. Then which of the following statements is(are) true ?
\begin{qoptions}
\opt{A}{There are infinitely many functions from S to T}
\opt{B}{There are infinitely many strictly increasing functions from S to T}
\opt{C}{The number of continuous functions from S to T is at most 786}
\opt{D}{Every continuous function from S to T is differentiable}
\end{qoptions}

\qn{6} Let $T_1$ and $T_2$ be two distinct common tangents to the ellipse $E:x^2+2y^2-2x-8y+3=0$ and the parabola $P:12x=y^2-4y+16$. Suppose that the tangent $T_1$ touches P and E at the points $A_1$ and $A_2$, respectively and the tangent $T_2$ touches P and E at the points $A_4$ and $A_3$, respectively. Then which of the following statements is(are) true ?
\begin{qoptions}
\opt{A}{The area of the quadrilateral $A_1A_2A_3A_4$ is 35 square units}
\opt{B}{The area of the quadrilateral $A_1A_2A_3A_4$ is 36 square units}
\opt{C}{The tangents $T_1$ and $T_2$ meet the x-axis at the point $(-2,2)$}
\opt{D}{The tangents $T_1$ and $T_2$ meet the x-axis at the point $(-5,2)$}
\end{qoptions}

\qn{7} Let $f:[0,1]\to[0,1]$ be the function defined by $36f(x)=12x^3-36x^2+20x+17$. Consider the square region $S=[0,1]\times[0,1]$. Let $G=\{(x,y)\in S:y>f(x)\}$ be called the green region and $R=\{(x,y)\in S:y<f(x)\}$ be called the red region. Let $L_h=\{(x,h)\in S:x\in[0,1]\}$ be the horizontal line drawn at a height $h\in[0,1]$. Then which of the following statements is(are) true ?
\begin{qoptions}
\opt{A}{There exists an $h\in\left[\dfrac14,\dfrac23\right]$ such that the area of the red region above the line $L_h$ equals the area of the green region below the line $L_h$}
\opt{B}{There exists an $h\in\left[\dfrac14,\dfrac23\right]$ such that the area of the red region above the line $L_h$ equals the area of the red region below the line $L_h$}
\opt{C}{There exists an $h\in\left[\dfrac14,\dfrac23\right]$ such that the area of the green region above the line $L_h$ equals the area of the red region below the line $L_h$}
\opt{D}{There exists an $h\in\left[\dfrac14,\dfrac23\right]$ such that the area of the green region above the line $L_h$ equals the area of the green region below the line $L_h$}
\end{qoptions}

\sect{Section-I(iii)}

\qn{8} Let $\alpha,\beta$ and $\gamma$ be real numbers. Consider the following system of linear equations\\
$x+2y+z=4$\\
$x+\alpha z=10$\\
$2x-3y+\beta z=\gamma+1$\\
Match each entry in \textbf{List - I} to the correct entries in \textbf{List-II}\par\nobreak
\begin{center}\small\setlength{\tabcolsep}{3pt}
\begin{tabular}{|c|p{3.6cm}|c|p{2.1cm}|}
\hline
\multicolumn{2}{|c|}{\textbf{List-I}} & \multicolumn{2}{c|}{\textbf{List-II}} \\
\hline
(P) & If $\beta=\dfrac12(7\alpha-3)$ and $\gamma=28$, then the system has & (1) & a unique solution \\
\hline
(Q) & If $\beta=\dfrac12(7\alpha-3)$ and $\gamma\neq28$, then the system has & (2) & no solution \\
\hline
(R) & If $\beta\neq\dfrac12(7\alpha-3)$ where $\alpha=1$ and $\gamma\neq28$, then the system has & (3) & infinitely many solutions \\
\hline
(S) & If $\beta\neq\dfrac12(7\alpha-3)$ where $\alpha=1$ and $\gamma=28$, then the system has & (4) & $x=10,\ y=-3$ and $z=0$ as a solution \\
\hline
 & & (5) & $x=-16,\ y=3$ and $z=0$ as a solution \\
\hline
\end{tabular}
\end{center}
The correct option is :
\begin{qoptions}
\opt{A}{P $\to$ 3; Q $\to$ 2; R $\to$ 5; S $\to$ 4}
\opt{B}{P $\to$ 3; Q $\to$ 2; R $\to$ 1; S $\to$ 4}
\opt{C}{P $\to$ 2; Q $\to$ 1; R $\to$ 4; S $\to$ 5}
\opt{D}{P $\to$ 2; Q $\to$ 1; R $\to$ 1; S $\to$ 3}
\end{qoptions}

\qn{9} Consider the given data with frequency distribution\\\par\nobreak
\begin{center}\small
\begin{tabular}{l|cccccc}
$x_i$ & 3 & 8 & 11 & 10 & 5 & 4 \\
$f_i$ & 5 & 2 & 3 & 2 & 4 & 4 \\
\end{tabular}
\end{center}
Match each entry in \textbf{List-I} to the correct entries in \textbf{List-II}.\par\nobreak
\begin{center}\small\setlength{\tabcolsep}{3pt}
\begin{tabular}{|c|p{3.7cm}|c|p{0.8cm}|}
\hline
 & \multicolumn{1}{c|}{\textbf{List-I}} & & \multicolumn{1}{c|}{\textbf{\footnotesize List-II}} \\
\hline
(P) & The mean of the above data is & (1) & 2.5 \\
\hline
(Q) & The median of the above data is & (2) & 5 \\
\hline
(R) & The mean deviation about the mean of the above data is & (3) & 6 \\
\hline
(S) & The mean deviation about the median of the above data is & (4) & 2.7 \\
\hline
 & & (5) & 2.4 \\
\hline
\end{tabular}
\end{center}
The correct option is :
\begin{qoptions}
\opt{A}{P $\to$ 3; Q $\to$ 3; R $\to$ 5; S $\to$ 5}
\opt{B}{P $\to$ 3; Q $\to$ 2; R $\to$ 1; S $\to$ 5}
\opt{C}{P $\to$ 2; Q $\to$ 3; R $\to$ 4; S $\to$ 1}
\opt{D}{P $\to$ 3; Q $\to$ 2; R $\to$ 4; S $\to$ 5}
\end{qoptions}

\qn{10} Let $\ell_1$ and $\ell_2$ be the lines $\vec r_1=\lambda\left(\hat i+\hat j+\hat k\right)$ and $\vec r_2=\left(\hat j-\hat k\right)+\mu\left(\hat i+\hat k\right)$, respectively. Let X be the set of all the planes H that contain the line $\ell_1$. For a plane H, let d(H) denote the smallest possible distance between the points of $\ell_2$ and H. Let $H_0$ be plane in X for which d($H_0$) is the maximum value of d(H) as H varies over all planes in X.\\
Match each entry in \textbf{List-I} to the correct entries in \textbf{List-II}.\par\nobreak
\begin{center}\small\setlength{\tabcolsep}{3pt}
\begin{tabular}{|c|p{3.9cm}|c|p{1.0cm}|}
\hline
 & \multicolumn{1}{c|}{\textbf{List-I}} & & \multicolumn{1}{c|}{\textbf{\footnotesize List-II}} \\
\hline
(P) & The value of d($H_0$) is & (1) & $\sqrt3$ \\
\hline
(Q) & The distance of the point (0,1,2) from $H_0$ is & (2) & $\dfrac1{\sqrt3}$ \\
\hline
(R) & The distance of origin from $H_0$ is & (3) & 0 \\
\hline
(S) & The distance of origin from the point of intersection of planes y = z, x = 1 and $H_0$ is & (4) & $\sqrt2$ \\
\hline
 & & (5) & $\dfrac1{\sqrt2}$ \\
\hline
\end{tabular}
\end{center}
The correct option is :
\begin{qoptions}
\opt{A}{P $\to$ 5; Q $\to$ 4; R $\to$ 3; S $\to$ 1}
\opt{B}{P $\to$ 2; Q $\to$ 4; R $\to$ 5; S $\to$ 1}
\opt{C}{P $\to$ 2; Q $\to$ 1; R $\to$ 3; S $\to$ 2}
\opt{D}{P $\to$ 5; Q $\to$ 1; R $\to$ 4; S $\to$ 2}
\end{qoptions}

\qn{11} Let z be complex number satisfying $27|z|^3+18z^2+12\bar z-8=0$, where $\bar z$ denotes the complex conjugate of z. Let the imaginary part of z be nonzero.\\
Match each entry in \textbf{List-I} to the correct entries in \textbf{List-II}.\par\nobreak
\begin{center}\small\setlength{\tabcolsep}{3pt}
\begin{tabular}{|c|p{3.0cm}|c|p{0.8cm}|}
\hline
 & \multicolumn{1}{c|}{\textbf{List-I}} & & \multicolumn{1}{c|}{\textbf{\footnotesize List-II}} \\
\hline
(P) & $9\,|z|^2$ is equal to & (1) & 12 \\
\hline
(Q) & $9\,|z-\bar z|^2$ is equal to & (2) & 4 \\
\hline
(R) & $9\left(|z|^2+|z+\bar z|^2\right)$ is equal to & (3) & 8 \\
\hline
(S) & $|3z+1|^2$ is equal to & (4) & 10 \\
\hline
 & & (5) & 7 \\
\hline
\end{tabular}
\end{center}
The correct option is :
\begin{qoptions}
\opt{A}{P $\to$ 1; Q $\to$ 3; R $\to$ 5; S $\to$ 4}
\opt{B}{P $\to$ 2; Q $\to$ 4; R $\to$ 5; S $\to$ 1}
\opt{C}{P $\to$ 2; Q $\to$ 1; R $\to$ 3; S $\to$ 5}
\opt{D}{P $\to$ 2; Q $\to$ 3; R $\to$ 5; S $\to$ 4}
\end{qoptions}

\sect{Section-II}

\qn{1} Let $\tan^{-1}(x)\in\left(-\dfrac\pi2,\dfrac\pi2\right)$, for $x\in\mathbb R$. Then the number of real solutions of the equation\\
$\sqrt{1+\cos(2x)}=\sqrt2\tan^{-1}\left(\tan\dfrac x6\right)$ in the set $(-9\pi,-3\pi)\cup(-3\pi,3\pi)\cup(3\pi,9\pi)$ is equal to

\qn{2} Let $n\geq2$ be a natural number and $f:[0,1]\to\mathbb R$ be the function defined by\\
$f(x)=\begin{cases}n(1-2nx)&\text{if }0\leqslant x\leqslant\frac1{2n}\\ 2n(2nx-1)&\text{if }\frac1{2n}\leqslant x\leqslant\frac3{4n}\\ 4n(1-nx)&\text{if }\frac3{4n}\leqslant x\leqslant\frac1n\\ \frac n{n-1}(nx-1)&\text{if }\frac1n\leqslant x\leqslant1\end{cases}$\\
If n is such that the area of the region bounded by the curves x = 0 , x = 1 , y = 0 and y = f (x) is 6, then the maximum value of the function f is

\qn{3} Let $\overbrace{75\dots57}^{r}$ denote the $(r+2)$ digit number where the first and the last digits are 7 and the remaining r digits are 5. Consider the sum $S=77+757+7557+\dots+\overbrace{75\dots57}^{100}$. If $S=\dfrac{\overbrace{75\dots57}^{101}+m}n$, where m and n are natural numbers less than 3000, then the value of $m+n$ is

\qn{4} Let $A=\left\{\dfrac{133-i(900\cos\theta+1686\sin\theta)}{7-3i\cos\theta}:\theta\in\mathbb R\right\}$. If A contains exactly one positive integer n, then the value of n is

\qn{5} Let P be the plane $\sqrt3x+2y+3z=16$ and let\\
$S=\left\{\alpha\hat i+\beta\hat j+\gamma\hat k:\alpha^2+\beta^2+\gamma^2=1\right.$\\ $\left.\text{and the distance of }(\alpha,\beta,\gamma)\text{ from the plane P is }\dfrac72\right\}$. Let $\vec u,\vec v$ and $\vec w$ be three distinct vectors in S such that $|\vec u-\vec v|=|\vec v-\vec w|=|\vec w-\vec u|$. Let V be the volume of the parallelepiped determined by vectors $\vec u,\vec v$ and $\vec w$. Then the value of $\dfrac{16}{\sqrt3}V$ is

\qn{6} Let a and b be two nonzero real numbers. If the coefficient of $x^7$ in the expansion of $\left(ax^3+\dfrac{70}{27bx^2}\right)^4$ is equal to the coefficient of $x^{-6}$ in the expansion of $\left(ax^2-\dfrac1{bx^3}\right)^7$, then the value of 2b is

\papersec{2}
\partsec{1}{PHYSICS}
\sect{Section-I(i)}

\qn{1} An electric dipole is formed by two charges $+q$ and $-q$ located in xy-plane at (0, 2) mm and (0, $-$2) mm, respectively, as shown in the figure. The electric potential at point P(100, 100) mm due to the dipole is $V_0$. The charges $+q$ and $-q$ are then moved to the points ($-$1, 0) mm and (1, 0) mm, respectively. What is the value of electric potential at P due to the new dipole ?
\qfig{p2_phy_s1_q01_1}
\begin{qoptions}
\opt{A}{$-\dfrac{V_0}4$}
\opt{B}{$\dfrac{V_0}2$}
\opt{C}{$\dfrac{V_0}{\sqrt2}$}
\opt{D}{$-\dfrac{V_0}2$}
\end{qoptions}

\qn{2} Density is expressed in terms of three derived quantities, namely, the gravitational constant G, Planck's constant h and the speed of light c, as $\rho=c^\alpha h^\beta G^\gamma$. Which of the following is the correct option?
\begin{qoptions}
\opt{A}{$\alpha=5,\ \beta=-1,\ \gamma=-2$}
\opt{B}{$\alpha=-5,\ \beta=-1,\ \gamma=-2$}
\opt{C}{$\alpha=5,\ \beta=-1,\ \gamma=2$}
\opt{D}{$\alpha=-5,\ \beta=1,\ \gamma=-2$}
\end{qoptions}

\qn{3} A particle of mass m is moving in the $xy$-plane such that its velocity at a point $(x,y)$ is given as $\vec v=\alpha(y\hat x+x\hat y)$ where $\alpha$ is a non-zero constant. What is the torque acting on the particle about the origin ?
\begin{qoptions}
\opt{A}{$\vec0$}
\opt{B}{$2m\alpha xy\hat k$}
\opt{C}{$m\alpha xy\hat k$}
\opt{D}{$-2m\alpha xy\hat k$}
\end{qoptions}

\qn{4} An ideal gas is in thermodynamic equilibrium. The number of degrees of freedom of a molecule of the gas is n. The internal energy of one mole of the gas is $U_n$ and the speed of sound in the gas is $v_n$. At a fixed temperature and pressure, which of the following is the correct option ?
\begin{qoptions}
\opt{A}{$v_3<v_6$ and $U_3>U_6$}
\opt{B}{$v_5>v_3$ and $U_3>U_5$}
\opt{C}{$v_5>v_7$ and $U_5<U_7$}
\opt{D}{$v_6<v_7$ and $U_6<U_7$}
\end{qoptions}

\sect{Section-I(ii)}

\qn{5} A monochromatic light wave is incident normally on a thin glass slab of thickness d, as shown in the figure. The refractive index of the slab increases linearly from $n_1$ to $n_2$ over the height h. Which of the following statement(s) is (are) true about the light wave emerging out of the slab?
\qfig{p2_phy_s2_q05_1}
\begin{qoptions}
\opt{A}{It will deflect up by an angle $\tan^{-1}\left[\dfrac{\left(n_2^2-n_1^2\right)d}{2h}\right]$}
\opt{B}{It will deflect up by an angle $\tan^{-1}\left[\dfrac{(n_2-n_1)d}h\right]$}
\opt{C}{It will not deflect.}
\opt{D}{The deflection angle depends only on $(n_2-n_1)$ and not on the individual values of $n_1$ and $n_2$.}
\end{qoptions}

\qn{6} An spherical shell of mass M, inner radius a and outer radius b is placed on a horizontal surface with coefficient of friction $\mu$, as shown in the figure. At some time, an impulse $J_0\hat x$ is applied at a height h above the center of the shell. If h = $h_m$ then the shell rolls without slipping along the x-axis. Which of the following statement(s) is(are) correct ?
\qfig{p2_phy_s2_q06_1}
\begin{qoptions}
\opt{A}{For $\mu\neq0$ and $a\to0$, $h_m=b/2$}
\opt{B}{For $\mu\neq0$ and $a\to b$, $h_m=2b/3$}
\opt{C}{For $h=h_m$ the initial angular velocity is inversely proportional to on the outer radius $b$.}
\opt{D}{For $\mu=0$ and $h=0$, the shell always slides without rolling.}
\end{qoptions}

\qn{7} The magnetic field associated with an electromagnetic wave propagating in a dielectric medium is given by $\vec B=10^{-7}\left(\hat x+\hat y\right)\sin\left[2\pi\left(5\times10^{14}t-\dfrac{10^7}3z\right)\right]$T. Which of the following option(s) is(are) correct? [Given: The speed of light in vacuum, $c=3\times10^8\ ms^{-1}$]
\begin{qoptions}
\opt{A}{$E_x=-15\sin\left[2\pi\left(5\times10^{14}t-\dfrac{10^7}3z\right)\right]$ v/m}
\opt{B}{$E_y=-15\sin\left[2\pi\left(5\times10^{14}t-\dfrac{10^7}3z\right)\right]$ v/m}
\opt{C}{The wave is polarized in the $xy$-plane with electric field at an angle of 45$^\circ$ with respect to the positive $x$-axis.}
\opt{D}{The refractive index of the medium is 2.}
\end{qoptions}

\sect{Section-II(i)}

\paraq{1}{2}
$S_1$ and $S_2$ are two identical sound sources of frequency 536 Hz. The source $S_1$ is located at $O$ and $S_2$ moves clockwise with a uniform speed 4 ms$^{-1}$ on a circular path of radius R around $O$, as shown in the figure. There are three points $P$, $Q$ and $R$ such that $P$ and $R$ are opposite at a distance of $\dfrac{3R}4$ from O while $Q$ is equidistant from them. A sound detector is placed at point $P$. The source $S_1$ can move along direction $OP$. [Given: The speed of sound in air is 324 ms$^{-1}$]
\qfig{p2_phy_s3_para1_2_1}

\qn{1} When only $S_2$ is emitting sound and it is at $Q$, the frequency of sound measured by the detector in Hz is \underline{\hspace{1cm}}.

\qn{2} Consider both sources emitting sound. When $S_2$ is above R and $S_1$ approaches the detector with a speed 4 ms$^{-1}$, the beat frequency measured by the detector is \underline{\hspace{1cm}}Hz.

\paraq{3}{4}
A cylindrical furnace has height ($H$) and diameter ($D$) both 1 m. It is maintained at temperature 360 K. The air gets heated inside the furnace at constant pressure $P_a$ and its temperature becomes $T=360$ K. The hot air with density $\rho$ rises up a vertical chimney of diameter $d=0.1$ m and height $h=9$ m above the furnace and exits the chimney (see the figure). As a result, atmospheric air of density $\rho_a=1.2$ kg m$^{-3}$, pressure $P_a$ and temperature $T_a=300$ K enters the furnace. Assume air as an ideal gas, neglect the variations in $\rho$ and $T$ inside the chimney and the furnace. Also ignore the viscous effects. Assume that just outside the furnace at the ground level pressure $P_a$ is atmospheric. [Given: The acceleration due to gravity $g=10$ ms$^{-2}$ and $\pi=3.14$]
\qfig{p2_phy_s3_para3_4_1}

\qn{3} Considering the air flow to be streamline, the steady mass flow rate of air exiting the chimney is \underline{\hspace{1cm}}gm s$^{-1}$.

\qn{4} When the chimney is closed using a cap at the top, a pressure difference $\Delta P$ develops between the top and the bottom surfaces of the cap. If the changes in the temperature and density of the hot air, due to the stoppage of air flow, are negligible then the value of $\Delta P$ is \underline{\hspace{1cm}} Nm$^{-2}$.

\sect{Section-II(ii)}

\qn{5} A thin circular coin of mass 5 gm and radius 4/3 cm is initially in a horizontal $xy$-plane. The coin is tossed vertically up (+$z$ direction) by applying an impulse of $\sqrt{\dfrac\pi2}\times10^{-2}$N $-$ s at a distance 1 cm from its center. The coin spins about its diameter and moves along the +$z$ direction. By the time the coin reaches back to its initial position, it completes $n$ rotations. The value of $n$ is :- [Given: The acceleration due to gravity g = 10 ms$^{-2}$]
\qfig{p2_phy_s4_q05_1}

\qn{6} A rectangular conducting loop of length 4 cm and width 2 cm is in the $xy$-plane, as shown in the figure. It is being moved away from a thin and long conducting wire along the direction $\dfrac45\hat x+\dfrac35\hat y$ with a constant speed $v$. The wire is carrying a steady current I = 10 A in the positive $x$-direction. A current of 10 $\mu$A flows through the loop when it is at a distance $d$ = 8 cm from the wire. If the resistance of the loop is 0.1 $\Omega$, then the value of $v$ is \underline{\hspace{1cm}} ms$^{-1}$. [Given: The permeability of free space $\mu_0=4\pi\times10^{-7}\ \text{NA}^{-2}$]
\qfig{p2_phy_s4_q06_1}

\qn{7} A string of length 1 m and mass x $\times$ $10^{-5}$ kg is under tension of 5N. When the string vibrates, two successive harmonics are found to occur at frequencies 750 Hz and 1000 Hz. The value of x is \underline{\hspace{1cm}}.

\qn{8} An incompressible liquid is kept in a container having a weightless piston with a hole. A capillary tube of inner radius 0.3 mm is dipped vertically into the liquid through the airtight piston hole, as shown in the figure. The air in the container is isothermally compressed from its original volume $V_0$ to $\dfrac{100}{101}V_0$ with the movable piston. Considering air as an ideal gas, the height ($h$) of the liquid column in the capillary above the liquid level in cm is \underline{\hspace{1cm}}.\\
\relax[Given: Surface tension of the liquid is 0.075 Nm$^{-1}$, atmospheric pressure is $10^5$ Nm$^{-2}$, acceleration due to gravity ($g$) is 10 ms$^{-2}$, density of the liquid is $10^3$ kgm$^{-3}$ and contact angle of capillary surface with the liquid is zero]
\qfig{p2_phy_s4_q08_1}

\qn{9} In a radioactive decay process, the activity is defined as $A=-\dfrac{dN}{dt}$, where $N(t)$ is the number of radioactive nuclei at time $t$. Two radioactive sources, $S_1$ and $S_2$ have same activity at time $t=0$. At a later time, the activities of $S_1$ and $S_2$ are $A_1$ and $A_2$, respectively. Half life of $S_1$ is 300 sec and half life of $S_2$ is 150 sec. At t = 600 sec, the ratio $A_1/A_2$ is \underline{\hspace{1cm}}.

\qn{10} Two mole of an ideal monatomic gas undergoes two different cyclic processes I and II, as shown in the $P$--$V$ diagrams below. In cycle I, processes $a$, $b$, $c$ and $d$ are isobaric, isothermal, isobaric and isochoric, respectively. In cycle II, processes $a'$, $b'$, $c'$ and $d'$ are isothermal, isochoric, isobaric and isochoric, respectively. The total heat exchange during cycle I is $H_1$ and that during cycle II is $H_2$. The ratio $\dfrac{H_1}{H_2}$ is \underline{\hspace{1cm}}.
\qfig{p2_phy_s4_q10_1}

\partsec{2}{CHEMISTRY}
\sect{Section-I(i)}

\qn{1} The correct molecular orbital diagram for O$_2$ molecule (if sp mixing is operative) in the ground state is
\begin{qoptions}
\opt{A}{\qfigi{p2_chem_s1_q01_A}}
\opt{B}{\qfigi{p2_chem_s1_q01_B}}
\opt{C}{\qfigi{p2_chem_s1_q01_C}}
\opt{D}{\qfigi{p2_chem_s1_q01_D}}
\end{qoptions}

\qn{2} Consider the following statements related to colloids.\\
(I) Lyophobic colloids can not be coagulated.\\
(II) For gel, both the dispersed phase and the dispersion medium are liquid.\\
(III) Micelles can not be formed below Kraft temperature.\\
(IV) Tyndall effect can be observed from a colloidal solution with dispersed phase having the different refractive index as that of the dispersion medium.\\
The option with the correct set of statements is
\begin{qoptions}
\opt{A}{(I) and (II)}
\opt{B}{(II) and (III)}
\opt{C}{(III) and (IV)}
\opt{D}{(II) and (IV)}
\end{qoptions}

\qn{3} In the following reactions, \textbf{W}, \textbf{X}, \textbf{Y}, and \textbf{Z} are the major products.
\qfig{p2_chem_s1_q03_1}
The incorrect statement about \textbf{W}, \textbf{X}, \textbf{Y}, and \textbf{Z} is
\begin{qoptions}
\opt{A}{\textbf{W} has six carbons and it undergoes Haloform reaction.}
\opt{B}{\textbf{X} undergoes Kolbe's electrolysis to give an eight-carbon product.}
\opt{C}{\textbf{Y} has six carbons and it undergoes Cannizzaro reaction.}
\opt{D}{\textbf{Z} is a secondary amine with six carbons.}
\end{qoptions}

\qn{4} A disaccharide \textbf{X} cannot be oxidised by Fehling solution. The acid hydrolysis of \textbf{X} leads to a laevorotatory solution. The disaccharide \textbf{X} is
\begin{qoptions}
\opt{A}{\qfigi{p2_chem_s1_q04_A}}
\opt{B}{\qfigi{p2_chem_s1_q04_B}}
\opt{C}{\qfigi{p2_chem_s1_q04_C}}
\opt{D}{\qfigi{p2_chem_s1_q04_D}}
\end{qoptions}

\sect{Section-I(ii)}

\qn{5} The complex(es), which can exhibit any type of isomerism shown by $[\text{Co(en)}_3]\text{Cl}_3$ , is(are)\\
\relax[en = H$_2$NCH$_2$CH$_2$NH$_2$] and Gly = NH$_2-$CH$_2-$COO$^-$ ]
\begin{qoptions}
\opt{A}{$[\text{Pt(en)(SCN)}_2]$}
\opt{B}{$[\text{Zn(gly)}_2]$}
\opt{C}{$[\text{Pt(NH}_3)_2\text{Cl}_4]$}
\opt{D}{$[\text{Cr(en)}_2\text{(H}_2\text{O)(SO}_4)]^+$}
\end{qoptions}

\qn{6} Atoms of metals x, y, and z form face-centred cubic (fcc) unit cell of edge length $L_x$, body-centred cubic (bcc) unit cell of edge length $L_y$, and simple cubic unit cell of edge length $L_z$, respectively.\\
If $r_z=\dfrac{\sqrt3}2r_y$; $r_y=\dfrac8{\sqrt3}r_x$; $M_z=\dfrac32M_y$ and $M_z=3M_x$, then the correct statement (s) is (are)\\
\relax[Given : $M_x$, $M_y$, and $M_z$ are molar masses of metals x, y, and z, respectively.\\
$r_x$, $r_y$, and $r_z$ are atomic radii of metals x, y, and z, respectively.]
\begin{qoptions}
\opt{A}{C.N. of x $>$ y $>$ z}
\opt{B}{$L_x>L_y$}
\opt{C}{Density of y $<$ Density of z}
\opt{D}{Density of x $<$ Density of y}
\end{qoptions}

\qn{7} In the following reactions, \textbf{P, Q, R,} and \textbf{S} are the major products.
\qfig{p2_chem_s2_q07_1}
The correct statement (s) about \textbf{P, Q, R,} and \textbf{S} is (are)
\begin{qoptions}
\opt{A}{Compounds \textbf{P} and \textbf{S} are the same.}
\opt{B}{\textbf{P, Q, R} and \textbf{S} are dicarboxylic acid.}
\opt{C}{Compounds \textbf{Q} and \textbf{R} are the same.}
\opt{D}{\textbf{R} does \textbf{not} undergo aldol condensation and \textbf{S} does \textbf{not} undergo Cannizzaro reaction.}
\end{qoptions}

\sect{Section-II(i)}

\paraq{1}{2}
The enthalpy versus temperature plot for phases $\alpha$ and $\beta$ at 1 bar pressure is given\\
$H_T$ is enthalpy of phases at temperatures T.
\qfig{p2_chem_s3_para1_2_1}
The transition temperature for $\alpha$ to $\beta$ phase change is 600 K and $C_{P,\beta}-C_{P,\alpha}=1$ J mol$^{-1}$ K$^{-1}$. Assume ($C_{P,\beta}-C_{P,\alpha}$) is independent of temperature in the range of 200 to 700 K. $C_{P,\alpha}$ and $C_{P,\beta}$ are heat capacities of $\alpha$ and $\beta$ phases, respectively.

\qn{1} The value of enthalpy change, $H_\beta-H_\alpha$ (in J mol$^{-1}$), at 300 K is \underline{\hspace{0.6cm}}.

\qn{2} The value of entropy change, $|S_\beta-S_\alpha|$ (in J mol$^{-1}$ K$^{-1}$), at 300 K is \underline{\hspace{0.6cm}}.\\
\relax[Use : ln 2 = 0.69\\
Given : $S_\beta-S_\alpha=0$ at 0 K]

\paraq{3}{4}
A trinitro compound, 1, 3,5 tris-(4-nitrophenyl) benzene, on complete reaction with an excess of Fe/HCl gives major product, which on treatment with an excess of NaNO$_2$/HCl at 0$^\circ$C provides \textbf{P} as the product. \textbf{P}, upon treatment with excess of H$_2$O at room temperature, gives the product \textbf{Q}. Bromination of \textbf{Q} in aqueous medium then treatment with Zn/$\Delta$ furnishes the product \textbf{R}. The compound \textbf{P} upon treatment with an excess of aniline under mild acidic conditions gives the product \textbf{S}.

\qn{3} The number of heteroatoms present in one molecule of \textbf{R} is \underline{\hspace{0.6cm}}.\\
\relax[Use: Molar mass (in g mol$^{-1}$): H = 1, C = 12, N = 14, O = 16, Br = 80, Cl = 35.5\\
Atoms other than C and H are considered as heteroatoms]

\qn{4} The total number of nitrogen atoms present in one molecule of \textbf{S} is \underline{\hspace{0.6cm}}.\\
\relax[Use: Molar mass (in g mol$^{-1}$): H = 1, C = 12, N = 14, O = 16, Br = 80, Cl = 35.5]

\sect{Section-II(ii)}

\qn{5} H$_2$S (3 moles) reacts completely with acidified aqueous potassium dichromate solution. In this reaction, the number of moles of water produced is \textbf{x}, and the number of moles of electrons involved is \textbf{y}. The value of (\textbf{x} + \textbf{y}) is \underline{\hspace{1cm}}.

\qn{6} Among $[\text{I}_3]^-$, $[\text{SiO}_4]^{4-}$, SO$_2$Cl$_2$, XeF$_2$, SF$_4$, ClF$_3$, Fe(CO)$_5$, XeO$_2$F$_2$, $[\text{PtCl}_4]^{2-}$, XeF$_4$, and IO$_2$F$_2^-$, the total number of species having $sp^3d$ hybridised central atom is \underline{\hspace{1cm}}.

\qn{7} Consider the following molecules : Br$_3$O$_8$, CrO$_5$, H$_5$P$_3$O$_{10}$, Na$_2$S$_5$, H$_2$S$_4$O$_6$, H$_2$S$_5$O$_6$, and C$_3$O$_2$, Mn$_2$(CO)$_{10}$\\
Count the number of atoms existing in their zero oxidation state in each molecule. Their sum is \underline{\hspace{1cm}}.

\qn{8} For He$^+$, a transition takes place from the orbit of radius 26.45 pm to the orbit of radius 105.8 pm. The change in wavelength associated with e$^-$ in transition is \underline{\hspace{1cm}} pm.\\
(Round off your answer to nearest integer)\\
\relax[\textbf{Use} : $\pi=3$, Bohr radius, a = 52.9 pm]

\qn{9} 50 mL of 0.2 molal urea solution (density = 1.012 g mL$^{-1}$ at 300 K) is mixed with 500 mL of a solution containing 0.06 g of urea. Both the solutions were prepared in the same solvent. Then ($\Delta T_b\times100$) is \underline{\hspace{1cm}}.\\
\relax[Use : Molar mass of urea = 60 g mol$^{-1}$; gas constant, R = 62 L Torr K$^{-1}$ mol$^{-1}$ ; Assume, $\Delta_{mix}H=0$, $\Delta_{mix}V=0$, $K_b=0.5$ K kg mol$^{-1}$]\\
$\Delta T_b$ = elevation in boiling point of final solution.

\qn{10} The reaction of 4-methyloct-1-ene (\textbf{P}, 2.52 g) with HBr in the presence of $(\text{C}_6\text{H}_5\text{CO})_2\text{O}_2$ gives two isomeric bromides in a 9 : 1 ratio, with combined yield of 50\%. Of these, the entire amount of the primary alkyl bromide was reacted with an appropriate amount of diethylamine followed by treatment with eq. K$_2$CO$_3$ to given a non-ionic product \textbf{S} in 100\% yield. The mass (in mg) of \textbf{S} obtained is \underline{\hspace{1cm}}.\\
\relax[Use molar mass (in g mol$^{-1}$) : H = 1, C = 12, N = 14, Br = 80]

\partsec{3}{MATHEMATICS}
\sect{Section-I(i)}

\qn{1} Let $f:[1,\infty)\to\mathbb R$ be a differentiable function such that $f(1)=\dfrac13$ and $3\displaystyle\int_1^xf(t)\,dt=xf(x)-\dfrac{x^3}3$, $x\in[1,\infty)$. Let $e$ denote the base of the natural logarithm. Then
\begin{qoptions}
\opt{A}{$f(x)=3$ has 2 real solution}
\opt{B}{$f(x)=-\dfrac12$ has 2 real solution}
\opt{C}{Range of $f(x)$ is $\left[-\dfrac12e^{-5/3},\infty\right)$}
\opt{D}{$f(e)=4e^2/3$}
\end{qoptions}

\qn{2} Consider an experiment of tossing a coin repeatedly until the outcomes of two consecutive tosses are same. If the probability of a random toss resulting in head is $\dfrac13$, then the probability that the experiment
\begin{qoptions}
\opt{A}{Ends with head given that it starts with tail $=\dfrac2{21}$}
\opt{B}{Ends with head given that it starts with head $=\dfrac1{14}$}
\opt{C}{Ends with head given that it starts with tail $=\dfrac1{21}$}
\opt{D}{Ends with head given that it starts with head $=\dfrac27$}
\end{qoptions}

\qn{3} Let $M=(a_{ij})$, $i,j\in\{1,2,3\}$, be the $3\times3$ matrix such that $a_{ij}=1$ if $j+1$ is divisible by i, otherwise $a_{ij}=0$. Then which of the following statements is (are) true ?
\begin{qoptions}
\opt{A}{M is invertible}
\opt{B}{There exists a nonzero column matrix $\begin{pmatrix}a_1\\a_2\\a_3\end{pmatrix}$ such that $M\begin{pmatrix}a_1\\a_2\\a_3\end{pmatrix}=\begin{pmatrix}-a_1\\-a_2\\-a_3\end{pmatrix}$}
\opt{C}{The set $\{X\in\mathbb R^3:MX=0\}$ is singleton where $\mathbf 0=\begin{pmatrix}0\\0\\0\end{pmatrix}$}
\opt{D}{The matrix $(M-2I)$ is invertible, where I is the $3\times3$ identity matrix}
\end{qoptions}

\qn{4} For $x\in\mathbb R$ let $\tan^{-1}(x)\in\left(-\dfrac\pi2,\dfrac\pi2\right)$. Then the minimum value of the function $f:\mathbb R\to\mathbb R$ defined by $f(x)=\displaystyle\int_0^{x\tan^{-1}x}\dfrac{e^{(t-\cos t)}}{1+t^{2025}}dt$ is
\begin{qoptions}
\opt{A}{$-1$}
\opt{B}{$-\dfrac\pi4$}
\opt{C}{0}
\opt{D}{1}
\end{qoptions}

\sect{Section-I(ii)}

\qn{5} Let $f:(0,1)\to\mathbb R$ be the function defined as $f(x)=[5x]\left(x-\dfrac15\right)^2\left(x-\dfrac25\right)^3\left(x-\dfrac35\right)$, where $[x]$ denotes the greatest integer less than or equal to x. Then which of the following statements is(are) true?
\begin{qoptions}
\opt{A}{The function f is discontinuous exactly at one point in $(0,1)$}
\opt{B}{There is exactly one point in $(0,1)$ at which the function f is continuous but \textbf{NOT} differentiable}
\opt{C}{The function f is differentiable in $\left(0,\dfrac12\right)$}
\opt{D}{Function f is differentible in $\left(\dfrac25,\dfrac45\right)$}
\end{qoptions}

\qn{6} Let S be the set of all twice differentiable functions f from $\mathbb R$ to $\mathbb R$ such that $\dfrac{d^2f}{dx^2}(x)>0$ for all $x\in(-1,1)$. For $f\in S$, let $X_f$ be the number of points $x\in(-1,1)$ for which $f(x)=\ell n(x+1)$. Then which of the following statements is(are) true?
\begin{qoptions}
\opt{A}{There exists a function $f\in S$ such that $X_f=3$}
\opt{B}{For every function $f\in S$, we have $X_f\leq2$}
\opt{C}{There exists a function $f\in S$ such that $X_f=2$}
\opt{D}{There does \textbf{NOT} exist any function f in S such that $X_f=1$}
\end{qoptions}

\qn{7} Let the position vectors of the points P,Q,R and S be $\vec a=\hat i+2\hat j-5\hat k$, $\vec b=3\hat i+6\hat j+3\hat k$, $\vec c=\dfrac{17}5\hat i+\dfrac{16}5\hat j+7\hat k$ and $\vec d=2\hat i+\hat j+\hat k$, respectively. Then which of the following statements is true?
\begin{qoptions}
\opt{A}{The points P,Q,R and S are coplanar}
\opt{B}{$\dfrac{\vec b+2\vec d}3$ is the position vector of a point which divides PR internally in the ratio 5 : 4}
\opt{C}{$\dfrac{\vec b+2\vec d}3$ is the position vector of a point which divides PR externally in the ratio 5 : 4}
\opt{D}{The square of the magnitude of the vector $\vec b\times\vec d$ is 99}
\end{qoptions}

\sect{Section-II(i)}

\paraq{1}{2}
Consider on obtuse angled triangle ABC in which the difference between the largest and the smallest angle is $\dfrac\pi2$ and whose sides are in arithmetic progression. Suppose that the vertices of this triangle lie on a circle of radius 1.

\qn{1} If area of triangle is $\Delta$. Then $64\Delta^2=$

\qn{2} Then the inradius of the triangle ABC is

\paraq{3}{4}
Consider the $6\times6$ square in the figure. Let $A_1,A_2,\dots,A_{49}$ be the points of intersections (dots in the picture) in some order. We say that $A_i$ and $A_j$ are friends if they are adjacent along a row or along a column. Assume that each point $A_i$ has an equal chance of being chosen.
\qfig{p2_math_s3_para3_4_1}

\qn{3} Let $p_i$ be the probability that a randomly chosen point has $i$ many friends, $i=0,1,2,3,4$. Let $X$ be a random variable such that for $i=0,1,2,3,4$, the probability $P(X=i)=p_i$. Then the value of $7E(X)$ is

\qn{4} Two distinct points are chosen randomly out of the points $A_1,A_2,\dots,A_{49}$. Let $p$ be the probability that they are not friends but have a friend in common . Then the value of $21p$ is

\sect{Section-II(ii)}

\qn{5} For $x\in\mathbb R$, let $y(x)$ be a solution of the differential equation $x\sin x\dfrac{dy}{dx}-(\sin x)y=xy\cos x-2x^3\sin^2x$ such that let\\
$f:(0,\pi)\to\mathbb R$\\
$f(x)=\dfrac{y(x)}{\sin x}$\\
$f(1)=2$\\
Then the maximum value of the function f(x) is-

\qn{6} Let X be the set of all five digit numbers formed using 01122333. For example, 33320 is in X while 02233 and 22233 are not in X . Suppose that each element of X has an equal chance of being chosen. Let p be the conditional probability that an element chosen at random is a multiple of 20 given that it is a multiple of 3. Then the value of 61p is equal to

\qn{7} Let $A_1,A_2,A_3,\dots\dots\dots,A_{16}$ be the vertices of a regular hexadecagon that lie on a circle of radius 2. Let P be a point on the circle and let $PA_i$ denote the distance between the points P and $A_i$ for i = 1,2,....,16. If P varies over the circle, then the maximum value of the product $PA_1\,PA_3.PA_5\cdot\cdot\cdot\cdot\cdot PA_{15}$, is

\qn{8} Let $R=\left\{\begin{bmatrix}a&b&c\\b&3&d\\0&4&0\end{bmatrix};a,b,c,d\in\{0,1,2,3,6,9\}\right\}$. Then number of invertible matrix in R is

\qn{9} Let $C_1$ be the circle of radius 1 with center at the origin. Let $C_2$ be the circle of radius r with center at the point A = (4,1) , where $1<r<3$. Two distinct common tangents PQ and ST of $C_1$ and $C_2$ are drawn. The tangent PQ touches $C_1$ at P and $C_2$ at Q. The tangent ST touches $C_1$ at S and $C_2$ at T. Mid points of the line segments PQ and ST are joined to form a line which meets the x-axis at a point B. If AB = $\sqrt5$, then the value of $r^2$ is

\qn{10} For any $y\in\mathbb R$ let $\cot^{-1}(y)\in(\pi,2\pi)$ and $\tan^{-1}(y)\in\left(\dfrac\pi2,\dfrac{3\pi}2\right)$. Then the number of solutions of the equation $\tan^{-1}\left(\dfrac{6y}{9-y^2}\right)+\cot^{-1}\left(\dfrac{9-y^2}{6y}\right)=\dfrac{8\pi}3$ for $0<|y|<3$, is equal to

\end{multicols}
\end{document}
